% !TeX root = ../main.tex
\clearpage

\section{Materials and Methods}

This thesis combines hardware development, a four week field expedition to Brazil in September 2025, and long term captive recordings in Berlin.
The quantitative Results use only the Berlin dataset.
The field recordings are described here because they motivated much of the hardware work and provide the natural habitat counterpart for future analyses; they are not analysed quantitatively in this thesis.

Custom multi electrode loggers, foldable square grids, flexible electrode lines, waterproof housings, and flight compatible battery packs were built before departure (Section~\ref{sec:recording_hardware}).
The first two weeks of the expedition were spent in Manaus at Reserva Ducke and the Instituto Nacional de Pesquisas da Amaz{\^o}nia (INPA).
At Ducke, grids and lines were deployed at four stream sites near known \textit{Electrophorus} activity, including nest associated locations.
At INPA, loggers were tested and calibrated in an enclosure with captive eels.
The second two weeks were spent near Macap{\'a} in the Floresta Nacional (FLONA) do Amap{\'a}, with continuous deployments along the Rio Falsino and Rio Araguari.
There, nine square grid (or tetrahedron) sites and seven line sites were instrumented in parallel, with batteries exchanged every few days (Figure~\ref{fig:flona_grids}).

\begin{figure}[htbp]
  \centering
  \begin{subfigure}[t]{0.48\textwidth}
    \includegraphics[width=\textwidth]{field/grid_deployed_site4.jpg}
    \subcaption{Surface floating $4 \times 4$ grid with bottle floats and logger housing at a Rio Falsino bank site.}
  \end{subfigure}
  \hfill
  \begin{subfigure}[t]{0.48\textwidth}
    \includegraphics[width=\textwidth]{field/grid_overview_site8.jpg}
    \subcaption{Submerged grid at a Rio Araguari site, viewed from a boat during deployment.}
  \end{subfigure}
  \caption{Electrode grids deployed in FLONA.}
  \label{fig:flona_grids}
\end{figure}

Table~\ref{tab:recording_inventory} summarises recording effort for the three datasets.
Berlin provides the multi year captive baseline used throughout the Results.
Ducke and FLONA together contribute about 11\,TiB of field WAV data from up to 16 concurrent loggers; those recordings remain for follow up work that applies the same classification and activity pipeline under natural conditions.

\begin{table}[htbp]
  \centering
  \small
  \caption{Recording inventory for the Berlin captive dataset and the September 2025 Brazil field campaigns.
  Logger hours are approximate, based on the number of stored five minute WAV segments.}
  \label{tab:recording_inventory}
  \begin{tabularx}{\textwidth}{llll>{\raggedright\arraybackslash}X>{\raggedright\arraybackslash}X}
    \toprule
    Dataset & Location & Period & Loggers & Arrays / sites & Effort (raw size) \\
    \midrule
    Berlin & MfN tank & Nov 2023--Mar 2026 & 1--2 & 16-ch line(s) & $\sim$3770\,h over 165\,d (13.7\,TiB) \\
    Ducke & Reserva Ducke & 1--2 Sep 2025 & 7 & 4 sites (grids + lines) & $\sim$236\,h (1.3\,TiB) \\
    FLONA & Falsino \& Araguari & 15--22 Sep 2025 & 16 & 9 grids/tetra + 7 lines & $\sim$1970\,h (9.9\,TiB) \\
    \bottomrule
  \end{tabularx}
\end{table}

\subsection{Recording hardware}
\label{sec:recording_hardware}

Recordings used custom multi electrode loggers based on Teensy microcontrollers with 16 channel amplifier boards.
Software and board designs are available at \url{https://github.com/janscience/TeeGrid}, \url{https://github.com/janscience/TeeRec}, and \url{https://github.com/janscience/Teensy_Amp}. 
Each device digitizes 16 electrode channels continuously to disk as multi channel WAV files.

Electrodes were wound by hand from stainless steel wire (three turns). 
Three array geometries were built: lines, square grids, and one tetrahedron for exploratory experiments. 

\textbf{Grids.} Square frames with 180\,cm side length carried 16 electrodes in a $4 \times 4$ layout at 60\,cm spacing.
Grid frames were assembled from PVC tubing: tubes form the edges, straight segments use sleeved plug-in PVC joints, and corners use PVC cross pieces with glued-on connectors.
This construction was foldable for transport while remaining robust and quick to deploy; the packed frame folds to 60\,cm edge length.
For assembly, tubes were laid out in the target configuration, electrodes slipped onto the tubes, and cable runs cut with slack for joints plus about 150\,cm from grid to housing.
Stretchy ropes (four on inner bars, two on cross bars) tension the frame.
Cables were soldered to the matching electrodes, sealed with waterproof heat shrink, and routed under spiral tubing along cable-bearing tubes.

\textbf{Lines.} Flexible electrode lines use a cable bundle wrapped in spiral tubing.
Flexibility reduces mechanical damage in the field and simplifies packing relative to rigid frames.
Total length is 15.6\,m: the first electrode sits 150\,cm from the housing end, and the remaining electrodes follow at 80\,cm spacing, totalling 16 electrodes.
For assembly, 16 cables were cut to length, bundled, and wrapped in spiral tubing with electrode positions marked.
At each site the matching cable was pulled out of the wrap, the electrode wound onto tubing and cable, then soldered.
Waterproof heat shrink sealed each joint.
Spiral wrap was applied in inter-electrode segments with about 20\,cm overlap at electrodes for stability.
Tube diameter decreased toward the free end as fewer cables remained.

The cable to housing transition used a threaded gland.
To ensure that the transition glans was water proof, heat shrink was applied over the 16 cable bundle.
Then the gland volume was filled with epoxy resin, avoiding air bubbles and epoxy leakage that would stiffen the free cables or interfere with their fit in the housing.


\subsubsection{Waterproofing and power}

Solder joints were sealed with waterproof heat shrink. 
Electronics sat in PVC housings. 
These housings were waterproofed with threaded end caps and O-rings. 
Before closing, O-rings were cleaned and silicone paste was used to insure seamless sealing. 
The outside to inside cable interface was potted in epoxy as above.

Battery packs were built from seven Lithium ion cells wired in parallel via nickel sheet on like poles (positive to positive, negative to negative).
A protection circuit was soldered on after corrosion protection spray. 
Contacts were kept from overheating during soldering by using minimal temperature and heat exposure time. 
Each pack was heat shrunk, capped with cut PVC plates, heat shrunk again, and labeled. 
Packs were stored at about 3.4\,V rest charge. 
Capacity is 24.5\,Ah per pack, which is within typical hand luggage flight limits, and lasts for about 60\,h or 2.5 days of logger runtime.

\subsection{Berlin captive recordings}

The data that builds the basis of the analysis conducted in this thesis were recorded in collaboration with Peter Bartsch at the Museum f{\"u}r Naturkunde (Natural History Museum Berlin).
Two adult eels, one male and one female, were housed in a dual-pool tank constructed from two pools of about 3\,m diameter joined by an opening of about 2\,m.
One pool was kept as a brighter region and the other as a darker region of the tank.
In the Berlin setting most recordings were obtained using a 16-electrode line with 25\,cm spacing, totalling 3.75\,m in length.
Overall, the Berlin dataset covers session starts on 165 calendar days from November 2023 to March 2026 (Table~\ref{tab:recording_inventory}), about 13.7\,TiB of raw WAV and around 500\,GiB of processed NIX/H5 data.

\subsection{Video and audio synchronization}

Where video was available, logger time and camera time were aligned with the custom open-source tool videosync. 
An LED on the logger blinks pseudo-randomly with a known schedule. 
videosync uses a video region of interest (ROI) to track the blinking LED on the logger.
Blink timing is recovered from logger metadata. 
Then both traces are cross correlated to estimate the temporal offset. 
Highest cross correlation value is used to align the two time bases.
Logger audio was band-pass filtered within 100 to 3000\,Hz, then stitched across overlapping WAV segments and muxed onto the video for later behavioral scoring.

\subsection{Data analysis}

Analysis of the Berlin recordings has two main stages.
First, EOD pulses are detected in the continuous WAV traces to obtain the time of every discharge and a short waveform cutout around it (Figure~\ref{fig:eod_detection}).
Second, those pulses are analysed by waveform morphology and by discharge rate, as detailed in the subsections below.

\begin{figure}[htbp]
  \centering
  \includegraphics[width=0.92\textwidth]{detection/eod_detection_excerpt.png}
  \caption{EOD pulse detection on a short Berlin recording excerpt (1\,s, one electrode channel).
  Red markers in (A) show eeltracker pulse times.
  (B) zooms on one detected discharge and shades the short waveform cutout stored for later shape analysis.}
  \label{fig:eod_detection}
\end{figure}

\subsubsection{Peak detection and EOD labeling}

Continuous recordings were processed with eeltracker.
The tool detects candidate pulses, groups them into multi-channel events with short waveform cutouts, and writes them to NIX/HDF5.
Downstream analyses keep only events keep only events that have been classified as EODs.

\subsubsection{EOD waveform classification}

Accepted EODs were further labelled by waveform morphology into three classes: normal, wide, and double pulses (Figure~\ref{fig:prototypes}).
Normal pulses are short single peaks while wide pulses are longer single peaks with a slower falling edge or shoulder. 
Double pulses show two peaks separated by a shallow trough (training labels also included clear shoulders toward a second peak).
A seed set was labelled by hand according to these shapes.
Waveforms were then reduced by RobustPCA (at most 20 components) and classified with a random forest trained on those labels.
Rare classes (wide, double) were accepted only if their predicted probability was at least 0.40.
The resulting markers (\texttt{is\_double\_peak}, \texttt{is\_wide\_pulse}, and \texttt{special\_pulse\_class}) provide the class assignments used throughout the Results.

Classified pulses were characterised by half-maximum width and by mean and median prototype waveforms per class; class prevalence is summarised over the dataset (Figures~\ref{fig:pulse_shape_distribution}--\ref{fig:half_width}).
To test whether the double shape is consistent with a short doublet of normal pulses, the median double prototype was fit as the sum of two median normal templates.
Six nested models (A--F) relaxed constraints on the inter-component delay $\Delta t$ and on the component widths (Figure~\ref{fig:double_model_hierarchy}).
Model E was chosen as the working description: it recovers the two-peak envelope and trough of the median double with free $t_1$ and $t_2$, a shared free width $w$, and free amplitudes $a_1$ and $a_2$.
Allowing independent widths in model F did not improve the fit meaningfully.
Labeled pulses were also examined in Robust PCA space (Figure~\ref{fig:pca_space}).

\subsubsection{Activity rates}

Pulse times are binned at minute, hour, day, month, month since recording start, and year.
Minute and day bins were discarded in the final analysis because they did not add information beyond the hour and month bins.
Rates are counts divided by recording effort in each bin, separately for all EODs, double pulses, and wide pulses.
The recording effort describes the total logger-on time (in seconds) that falls into that bin, not calendar time: a month with only a few recording days contributes only those recorded seconds.
This converts raw pulse counts into rates in Hz and avoids treating more heavily sampled periods as more active.
A second electrode line was added in November 2025 to span bright and dark regions of the tank.
Dual-line recordings from 25 November 2025 onward enter rate analyses with pulse counts weighted by 0.5, so the second line does not double-count activity.
The multi-timescale panels (Figure~\ref{fig:activity_all_shapes}) are active-session medians of those per-session rates, with bootstrap 95\% CIs, not a single pooled rate per bin.
Here, an active session is one with a positive pulse rate in that bin; sessions with no recording effort or zero pulses are excluded from the median.
Yearly panels therefore report session-median rates in Hz, so partial calendar years 2023 and 2026 are not scaled by the fraction of the year recorded; those years remain samples of the months that were actually logged.

\subsubsection{Behavioral and environmental correlations}

% Special pulse rates are cross correlated with overall EOD rate.
% Mating related notes from session logs are aligned to recording time and compared with all pulse and double pulse rates.
Temperature and conductivity from tank water logs are correlated with daily and monthly rates.
% High voltage volleys are scored as rapid trains: primary rule $\geq 5$ waveform sub peaks with inter event interval $\leq 2$\,ms, with a fallback of $\geq 3$ pulse centers at $\leq 5$\,ms.
% Bright versus dark tank regions tagged in filenames are compared for shape fractions.
Feeding windows ($\pm$5\,min around logged feeding times) are contrasted with baseline minutes ($>$5\,min from feeding) using mean pulse rates and a two-sided Mann--Whitney $U$ test on the minute-level rates.
Peri-feeding trajectories are computed as well.
Mean feeding clock time is the circular mean of the first logged feeding-related timestamp per session, with a bootstrap 95\% CI.

\subsubsection{Position along the electrode line}

For each accepted pulse, head position along the Berlin line is taken as the electrode with the largest positive peak amplitude in the multi-channel snippet (25\,cm bins).
Channel 0, closest to the logger, lay in the bright pool; later electrodes extend toward the dark pool.
The bright versus dark boundary used for occupancy analyses sits at about 2.5\,m along this line.
Occupancy histograms and mean position are computed on the same timescales as the activity rates, including bright versus dark occupancy relative to that boundary.
Collapsed occupancy rates are electrode pulse counts divided by total recording seconds (the sum of day of year recording time bins).

\subsubsection{Code availability}

Hardware repositories are linked above. 
Peak detection and video sync live in eeltracker (\url{https://github.com/weygoldt/eeltracker}) and videosync (\url{https://github.com/weygoldt/videosync}). 
The Berlin analysis pipeline used for this thesis (special pulse classification through position plots, matching \texttt{run\_analysis.sh}) is at \url{https://github.com/saraheisele/mscthesis}.
